\documentclass{article}
\usepackage{PRIMEarxiv} % Core package for the PrimeAI Template
\usepackage{amsmath, amssymb, amsthm, bm, dcolumn} % Add amsthm here for the proof environment
\usepackage[numbers,sort&compress]{natbib} % Natbib for citations
\usepackage{graphicx} % For high-quality images
\usepackage{multicol} % Multi-column figures/tables
\usepackage{hyperref} % Hyperlinks
\usepackage{listings} % Code listings
\usepackage{authblk} % For structured affiliations
% Define theorem style (optional)
\theoremstyle{plain}
\newtheorem{theorem}{Theorem}
\renewcommand{\qedsymbol}{}

% Custom commands for primes and qubit encoding
\newcommand{\primeQubit}[1]{|p_{#1}\rangle}
\newcommand{\primeGate}[1]{U_{p_{#1}}}

\usepackage{wrapfig}
\usepackage[pscoord]{eso-pic}
\usepackage[fulladjust]{marginnote}
\reversemarginpar

% Typesetting improvements without footnote patching
\usepackage[protrusion=true, expansion=true, tracking=false]{microtype}
\microtypecontext{spacing=nonfrench}

% Line numbers
\usepackage[right]{lineno}

% Text layout - adjust as needed
\raggedright
\setlength{\parindent}{0.5cm}
\textwidth 5.25in 
\textheight 8.75in

% Set double spacing
\usepackage{setspace} 
\doublespacing

% Adjust width for specific content
\usepackage{changepage}

% Adjust caption style
\usepackage[aboveskip=1pt,labelfont=bf,labelsep=period,singlelinecheck=off]{caption}

% Remove brackets from references
\makeatletter
\renewcommand{\@biblabel}[1]{\quad#1.}
\makeatother

% Header, footer, and page numbers
\usepackage{lastpage,fancyhdr}
\pagestyle{fancy}
\fancyhf{}
\fancyhead[L]{Preprint - PrimeAI Enhanced Template}
\fancyfoot[C]{\scriptsize Multiplicity Theory © 2024 Citizen Gardens \\ Licensed Under MIT and CC BY-NC-SA 4.0.}
\fancyfoot[R]{Page \thepage\ of \pageref{LastPage}}
\renewcommand{\footrule}{\hrule height 2pt \vspace{2mm}}

\begin{document}
\title{Unification of Tensor Frameworks in Multiplicity Theory}

\author{Citizen Gardens - The Foundation of Multiplicity}
\affil{info@citizengardens.org}

\maketitle

\begin{abstract}
This paper presents a unified tensor framework integrating Alena Tensors, Tensor Principal Component Analysis (PCA), and the Universal Multiplicity Equation (UME). By leveraging prime-based encoding, recursive tensor feedback, and quantum-inspired optimization, we establish a scalable methodology for high-dimensional tensor analysis. Applications include quantum computing, AI-driven optimization, and astrophysical simulations.
\end{abstract}

\section{Introduction}
Tensor methodologies are foundational in high-dimensional computing, enabling efficient representation of multi-modal data. Recent advances in Multiplicity Theory introduce prime-based encoding, tensor networks, and recursive feedback loops, significantly enhancing computational precision and scalability.

\subsection{Prime-Based Tensor Encoding}
A prime-based encoding method assigns each tensor element a unique prime representation:
\begin{equation}
\phi(T_{i_1,\dots,i_p}) = \prod_{k=1}^{p} p_{i_k},
\end{equation}
where $p_{i_k}$ are distinct primes ensuring modular and lossless representation.

\subsection{Universal Multiplicity Equation with Tensor Operations}
The Universal Multiplicity Equation extends tensor operations:
\begin{equation}
\frac{\partial \Psi(t)}{\partial t} = A(t) \Psi(t) + B(t) \int_{0}^{t} I(\tau) d\tau + T (t) \otimes \Psi(t) \otimes \Psi(t) + Q(t) \nabla^2 \Psi(t) + E(t).
\end{equation}

\subsection{Tensor Principal Component Analysis (PCA)}
Tensor PCA decomposes noisy high-dimensional data via spectral decomposition:
\begin{equation}
M (t+1) = f(M (t), R(t)), \quad R(t) = \phi(G) \cdot M (t).
\end{equation}

\section{Applications}
\subsection{Quantum Computing}
By integrating Alena tensors with Multiplicity Theory, we achieve enhanced quantum error correction using:
\begin{equation}
|\psi(\gamma, \beta)\rangle = U(C, \gamma)U(B, \beta)|\psi_0\rangle.
\end{equation}

\subsection{Artificial Intelligence and Machine Learning}
Recursive tensor feedback refines neural networks and AI models:
\begin{equation}
M(t+1) = f(M(t), R(t)) + \alpha T(M(t)).
\end{equation}

\subsection{Astrophysical Simulations}
Tensor networks enable large-scale cosmological modeling:
\begin{equation}
T_{i_1,\dots,i_p} = \sum_{k} \phi(p_{i_k}) \cdot \phi(G_k).
\end{equation}

\section{Conclusion}
The integration of Alena Tensors, Tensor PCA, and the Universal Multiplicity Equation into a singular framework advances computational efficiency in quantum computing, AI, and astrophysical modeling. Future research will explore prime-indexed quantum gates and recursive entanglement structures.

\bibliographystyle{plain}
\bibliography{references}
\end{document}
